\emph{(i)} By the series identity for $e^{B} x e^{-B}$ (\claimref{std-algebra.bivector-exponential-in-spin}(ii)), $\Ad_{e^{B}} = \exp(X)$ on $V$ with $X = \operatorname{ad}_{B}|_{V}$, which is a real linear map of $V$ because the commutator of a bivector with a vector is a vector (\claimref{std-algebra.bivector-lie-algebra}). Over $\C$, $X = P T P^{-1}$ with $T$ upper triangular with diagonal entries $\mu_1, \dots, \mu_n$ (the eigenvalues of $X$ with algebraic multiplicity), so $\exp(X) = P \exp(T) P^{-1}$ with $\exp(T)$ upper triangular with diagonal entries $e^{\mu_i}$: the eigenvalues of $\exp(X)$ with algebraic multiplicity are the $e^{\mu_i}$. Let $\lambda < 0$. If $e^{\mu} = \lambda$ then $\operatorname{Im}\mu$ is an odd multiple of $\pi$, so $\mu \ne \bar\mu$; since $X$ is real its characteristic polynomial has real coefficients, so $\bar\mu$ is an eigenvalue with the same multiplicity as $\mu$, and $e^{\bar\mu} = \overline{\lambda} = \lambda$. The multiset $\{\mu_i : e^{\mu_i} = \lambda\}$ is therefore closed under the fixed-point-free involution $\mu \mapsto \bar\mu$ and has even size. Finally $\Ad_{-e^{B}} = \Ad_{e^{B}}$ because $-1$ is central.

\emph{(ii)} $R_0 = \tfrac54 + \tfrac34 e_1e_3 = e_1\,(\tfrac54 e_1 + \tfrac34 e_3)$ with $Q(\tfrac54 e_1 + \tfrac34 e_3) = \tfrac{25}{16} - \tfrac{9}{16} = 1$, and $-e_1e_2e_3e_4 = (-e_1)\,e_2\,e_3\,e_4$; so $R$ is a product of six vectors with $Q = 1, 1, -1, -1, 1, 1$, and $R\tilde R = \prod Q(u_i) = 1$: $R$ is a rotor (\claimref{obj.std-algebra.rotor}). Since $(e_1e_3)^{2} = -Q(e_1)Q(e_3) = 1$, $R_0 = \cosh a + \sinh a\, e_1e_3$ with $e^{a} = \tfrac54 + \tfrac34 = 2$ (\claimref{std-algebra.bivector-exponential}). Write $\nu_{\pm} = e_1 \pm e_3$. As $e_1 e_3$ anticommutes with $e_1$ and $e_3$ and commutes with every vector orthogonal to both (\claimref{std-algebra.orthogonal-anticommute}), $R_0 x \tilde R_0 = R_0^{2} x$ on $\operatorname{span}(e_1, e_3)$ and $= x$ on its orthogonal complement; $R_0^{2} = \tfrac{17}{8} + \tfrac{15}{8} e_1e_3$, and $e_1e_3\,\nu_{\pm} = \mp\nu_{\pm}$ ($e_1e_3e_1 = -e_3 Q(e_1) = -e_3$, $e_1 e_3 e_3 = Q(e_3) e_1 = -e_1$), so $\Ad_{R_0}\nu_{\pm} = (\tfrac{17}{8} \mp \tfrac{15}{8})\nu_{\pm}$: the eigenvalues $\tfrac14$ on $\nu_{+}$ and $4$ on $\nu_{-}$. Let $\omega_4 = e_1e_2e_3e_4$ (the volume element of the $(2,2)$ block, of square $+1$; not the field theory's unit $I_4$). Each $e_j$, $j \le 4$, anticommutes with the three other factors, so $e_j \omega_4 = -\omega_4 e_j$, while a vector orthogonal to all four anticommutes with each factor and commutes with $\omega_4$; hence $\Ad_{\omega_4} = -1$ on $\operatorname{span}(e_1, \dots, e_4)$ and $+1$ on its complement, and $\omega_4$ commutes with the even element $R_0$. So $\Ad_R = \Ad_{\omega_4}\Ad_{R_0}$ has eigenvalue $-4$ on $\nu_{-}$, $-\tfrac14$ on $\nu_{+}$, $-1$ on $e_2$ and $e_4$, and $1$ on the remaining $n - 4$ directions; its characteristic polynomial is $(t + 4)(t + \tfrac14)(t + 1)^{2}(t - 1)^{n - 4}$, and $-4$ is a simple eigenvalue. If $R = \pm e^{B}$ for a bivector $B$, then $\Ad_R = \Ad_{e^{B}}$ would have $-4$ with even multiplicity by (i), a contradiction.

\emph{(iii)} For $n = 1$, $\Cl^{2} = 0$ so $e^{B} = 1$ for the only bivector $B = 0$, and $-1 = e_1\,(-Q(e_1)\,e_1)$ is a product of two vectors with $R\tilde R = Q(e_1)^{4} = 1$ (in $\Cl(0,1)$ this is $e_1 e_1$; note $e_1(-e_1) = +1$ there). In $\Cl(1,1)$ the even part is $\operatorname{span}\{1, e_1e_2\}$ with $(e_1e_2)^{2} = -Q(e_1)Q(e_2) = 1$ and $(e_1e_2)^{\sim} = -e_1e_2$, so $(a + b e_1e_2)(a + b e_1e_2)^{\sim} = a^{2} - b^{2}$; the bivectors are $\lambda e_1 e_2$ and $e^{\lambda e_1e_2} = \cosh\lambda + \sinh\lambda\, e_1e_2$ (\claimref{std-algebra.bivector-exponential}), whose scalar part is at least $1$. The rotor $-1 = e_1 (-e_1)$ has scalar part $-1$, so it is not an exponential; every rotor $a + b e_1e_2$ has $\lvert a \rvert \ge 1$ and is $\pm e^{\lambda e_1 e_2}$ with $\tanh\lambda = b/a$. In both cases every rotor is $\pm e^{B}$.

The check verifies the rotor of (ii), its six-vector factorisation, $R\tilde R = 1$ and the exact characteristic polynomial of $\Ad_R$ with the simple root $-4$ in every signature with $p, q \ge 2$ and $n \le 6$; the even multiplicity of the eigenvalue $-1$ for the exponential rotors $e^{\frac{\pi}{2}e_ie_j} = e_ie_j$ ($Q(e_i)Q(e_j) = 1$) in the same signatures; and (iii) in $\Cl(1,0)$, $\Cl(0,1)$ and $\Cl(1,1)$ at rational points of the hyperbola. The series identity behind (i) is checked with \claimref{std-algebra.bivector-exponential-in-spin}.

\emph{Remark (not part of the claim).} Only one direction is proved: $\min(p, q) \ge 2$ suffices for a rotor that is not $\pm e^{B}$. For $\min(p, q) \le 1$, (iii) and the Lorentzian case \claimref{std-algebra.bivector-exponential-in-spin}(iv) are consistent with every rotor being $\pm e^{B}$; in general this would follow from surjectivity of the exponential map onto $\mathrm{SO}_0(p, q)$ for $\min(p, q) \le 1$ (Nishikawa 1983, cited from memory) lifted through the kernel $\{\pm 1\}$.
